%! Constructing a Confidence Interval for a Population Proportion — Unit 6, Lesson 6.2 — Group Activity
\documentclass[10pt]{article}
\usepackage{apstats-article}
\usepackage{apstats-boxes}
\newcommand{\TallMath}[1]{$\displaystyle #1\rule[-1.4em]{0pt}{3.2em}$}

\begin{document}

\pageheader{Unit 6, Lesson 6.2}{Group Activity}
\namepartnerperiod

\vspace{0.15in}

\textit{For each scenario, complete the four-step inference procedure:
\textbf{State} the parameter, \textbf{Plan} (verify conditions), \textbf{Do} (compute the interval),
and \textbf{Conclude} (interpret in context).}

\begin{scenariobox}[Scenario 1 --- Teen Screen Time]{sky}
A researcher randomly selects 320 teenagers and asks whether they spend more than 4 hours
per day on screens. 176 say yes. Construct a 95\% confidence interval for the true proportion.

\textbf{State:} $p = $ \writelines{1}

\textbf{Plan:}
\begin{itemize}
  \item Random: \blank{7cm}
  \item Large Counts: $n\hat p = $ \blank{2.5cm}, $n(1-\hat p) = $ \blank{2.5cm} \quad Both $\ge 10$? \blank{1cm}
  \item 10\% rule: \blank{5cm}
\end{itemize}

\textbf{Do:} $\hat p = $ \blank{2cm}

\[
  \hat p \pm z^* \sqrt{\frac{\hat p(1-\hat p)}{n}}
  = \blank{1.5cm} \pm \blank{1.5cm}\sqrt{\frac{\blank{2cm}}{\blank{1.5cm}}}
  = \blank{1.5cm} \pm \blank{1.5cm}
  = (\blank{2.5cm},\; \blank{2.5cm})
\]

\textbf{Conclude:} \writelines{2}
\end{scenariobox}

\begin{scenariobox}[Scenario 2 --- Organic Food Preferences]{sky}
A random survey of 500 grocery shoppers finds 215 regularly buy organic produce. Construct
a 99\% confidence interval for the true proportion of shoppers who buy organic produce.

\textbf{State:} \writelines{1}

\textbf{Plan:} (verify all three conditions)
\begin{itemize}
  \item[] \writelines{3}
\end{itemize}

\textbf{Do:}
\[
  CI = \blank{1.5cm} \pm \blank{1.5cm}\sqrt{\frac{\blank{2cm}}{\blank{1.5cm}}}
     = (\blank{2.5cm},\; \blank{2.5cm})
\]

\textbf{Conclude:} \writelines{2}
\end{scenariobox}

\begin{scenariobox}[Scenario 3 --- Sample Size Planning]{sky}
A university wants to estimate the proportion of students who commute to campus. The
administration wants a 95\% confidence interval with a margin of error no larger than 0.04.
No prior estimate is available.

\begin{enumerate}[label=\alph*.]
  \item What value should you use for $\hat p$ when no prior estimate exists, and why?
        \writelines{2}

  \item Calculate the minimum required sample size.
        \[n \ge \left(\frac{z^*}{m}\right)^2 \hat p(1-\hat p) = \left(\frac{\blank{1.5cm}}{\blank{1.5cm}}\right)^2 (\blank{1.5cm})(\blank{1.5cm}) = \blank{3cm}\]
        Minimum $n = $ \blank{2cm} (round \underline{\hspace{1.5cm}}).
\end{enumerate}
\end{scenariobox}

\begin{reflectionbox}
Compare your intervals from Scenarios 1 and 2. Even though Scenario 2 used a higher
confidence level, how did the width compare? What two competing forces explain this?
\writelines{3}
\end{reflectionbox}

\end{document}
